% =============================================================================
% exercises/ejercicios_cap01a.tex
% Ejercicios propuestos — Capítulo 1, Sección A
% Temas: Leyes de los gases, ley general, peso molecular, conversiones
% =============================================================================

\section*{Ejercicios --- Secci\'on A: Gases y conversiones}
\addcontentsline{toc}{section}{Ejercicios: Gases y conversiones}
\label{sec:ejercicios_cap01a}

\begin{bloque_ejercicios}[Ejercicios A --- Gases y ecuaciones de estado]

\begin{ejercicios}

% ----- Ejercicio A1: Ley de Boyle -------------------------------------------
\item Un gas se comprime isot\'ermicamente a \SI{25}{\celsius} desde
$P_1 = \SI{0.8}{\atm}$ y $V_1 = \SI{15}{\liter}$ hasta una presi\'on
final $P_2 = \SI{2.4}{\atm}$.

\begin{enumerate}[label=\alph*)]
    \item Identifique qu\'e ley aplica e indique qu\'e variable
          permanece constante.
    \item Calcule el volumen final $V_2$.
    \item Si se comprime adicionalmente el gas hasta $V_3 = \SI{2}{\liter}$,
          ?`cu\'al ser\'ia la presi\'on final?
\end{enumerate}

\textit{[Respuesta: (b) $V_2 = P_1V_1/P_2 = 0.8\times15/2.4
= \SI{5}{\liter}$;
(c) $P_3 = P_1V_1/V_3 = 0.8\times15/2 = \SI{6}{\atm}$]}

% ----- Ejercicio A2: Ley de Charles -----------------------------------------
\item Un gas a presi\'on constante de \SI{1}{\atm} tiene un volumen de
\SI{250}{\milli\liter} a \SI{27}{\celsius}.

\begin{enumerate}[label=\alph*)]
    \item Convierta \SI{27}{\celsius} a Kelvin.
    \item Calcule el volumen que ocupar\'a el gas si se calienta a
          \SI{127}{\celsius}.
    \item ?`A qu\'e temperatura (en \si{\celsius}) ocupar\'ia el gas
          un volumen de \SI{375}{\milli\liter}?
\end{enumerate}

\textit{[Respuesta:
(a) \SI{300}{\kelvin};
(b) $V_2 = 250\times400/300 = \SI{333.3}{\milli\liter}$;
(c) $T_2 = 375\times300/250 = \SI{450}{\kelvin} = \SI{177}{\celsius}$]}

% ----- Ejercicio A3: Ley combinada de gases ---------------------------------
\item Un gas ocupa un volumen de \SI{3.5}{\liter} a una temperatura de
\SI{20}{\celsius} y una presi\'on de \SI{740}{\mmHg}. Se le somete
simult\'aneamente a una temperatura de \SI{80}{\celsius} y una presi\'on
de \SI{1}{\atm}.

\begin{enumerate}[label=\alph*)]
    \item Escriba la ley combinada de los gases e identifique qu\'e
          variables cambian.
    \item Exprese las presiones en las mismas unidades y calcule
          el volumen final $V_2$.
\end{enumerate}

\textit{[Respuesta:
(b) $P_1 = 740/760\,\si{\atm} = 0.9737\,\si{\atm}$;
$V_2 = (0.9737\times3.5\times353.15)/(1\times293.15)
= \SI{4.10}{\liter}$]}

% ----- Ejercicio A4: Ley general PV = nRT -----------------------------------
\item Utilizando la ley general de los gases $PV = nRT$
($R = 0.082057\ \si{\liter\atm\per\mole\per\kelvin}$):

\begin{enumerate}[label=\alph*)]
    \item Calcule el volumen que ocupa \SI{1}{\mole} de gas ideal
          a \SI{0}{\celsius} y \SI{1}{\atm} (condiciones normales, CN).
    \item Calcule el volumen que ocupa \SI{1}{\mole} de gas ideal
          a \SI{25}{\celsius} y \SI{1}{\atm} (condiciones est\'andar, CE).
    \item Calcule el volumen que ocupan \SI{44}{\gram} de \ce{CO2}
          a la presi\'on de la Ciudad de M\'exico ($P = \SI{0.77}{\atm}$)
          y \SI{20}{\celsius}.
    \item Un recipiente r\'igido de \SI{5}{\liter} contiene
          \SI{10}{\gram} de \ce{N2} a \SI{25}{\celsius}.
          Calcule la presi\'on interna en atm\'osferas.
\end{enumerate}

\textit{[Respuesta:
(a) $V = \SI{22.41}{\liter}$;
(b) $V = \SI{24.47}{\liter}$;
(c) $n = 44/44 = \SI{1}{\mole}$;
    $V = 1\times0.082057\times293.15/0.77 = \SI{31.2}{\liter}$;
(d) $n = 10/28 = 0.357\,\si{\mole}$;
    $P = 0.357\times0.082057\times298.15/5 = \SI{1.75}{\atm}$]}

% ----- Ejercicio A5: Densidad de un gas -------------------------------------
\item Con la ecuaci\'on de densidad del gas $\rho = PM/RT$:

\begin{enumerate}[label=\alph*)]
    \item Calcule la densidad del aire a \SI{20}{\celsius} y
          \SI{1}{\atm}, considerando que el peso molecular promedio
          del aire es \SI{29}{\gram\per\mole}.
    \item ?`Cu\'al ser\'ia la densidad del mismo aire a
          \SI{100}{\celsius} y \SI{1}{\atm}?
    \item Un gas desconocido tiene una densidad de
          \SI{1.964}{\kilo\gram\per\cubic\metre} a \SI{0}{\celsius}
          y \SI{1}{\atm}. Calcule su peso molecular e identifique
          de qu\'e gas se trata.
\end{enumerate}

\textit{[Respuesta:
(a) $\rho = 1\times29/(0.082057\times293.15)
= \SI{1.205}{\kilo\gram\per\cubic\metre}$;
(b) $\rho = 1\times29/(0.082057\times373.15)
= \SI{0.946}{\kilo\gram\per\cubic\metre}$;
(c) $M = \rho RT/P
= 1.964\times0.082057\times273.15/1
\approx \SI{44}{\gram\per\mole}$\ $\Rightarrow$ \ce{CO2}]}

% ----- Ejercicio A6: Peso molecular promedio --------------------------------
\item Un gas de combusti\'on tiene la siguiente composici\'on
en base seca (por ciento en volumen):
12\,\% \ce{CO2}, 8\,\% \ce{O2}, 6\,\% \ce{CO} y 74\,\% \ce{N2}.

\begin{enumerate}[label=\alph*)]
    \item Calcule la fracci\'on mol de cada componente
          usando $y_i = \%_{vi}/100$.
    \item Calcule el peso molecular promedio base seca
          $\overline{M}_s = \sum y_i M_i$.
    \item Si se sabe que la humedad del gas es del 10\,\%,
          calcule el peso molecular base h\'umeda
          $\overline{M}_w$ usando la \textbf{Ec.~\ref{hum}}.
    \item Calcule la densidad del gas h\'umedo a \SI{150}{\celsius}
          y \SI{1}{\atm}.
\end{enumerate}

\textit{[Respuesta parcial:
(b) $\overline{M}_s = 0.12(44)+0.08(32)+0.06(28)+0.74(28)
= 5.28+2.56+1.68+20.72 = \SI{30.24}{\gram\per\mole}$;
(c) $\overline{M}_w = 0.10(18)+0.90(30.24)
= \SI{29.02}{\gram\per\mole}$]}

% ----- Ejercicio A7: Número de moléculas y Avogadro -------------------------
\item La concentraci\'on de \ce{O3} en un d\'ia de contaminaci\'on
severa es de \SI{0.120}{ppm_v} en la Ciudad de M\'exico
($P = \SI{0.77}{\atm}$, $T = \SI{25}{\celsius}$).

\begin{enumerate}[label=\alph*)]
    \item Convierta la concentraci\'on de \ce{O3} de ppm$_v$ a
          \si{\milli\gram\per\cubic\metre} usando el factor
          $\text{ppm}_v \times M/24.45$ (a \SI{25}{\celsius}
          y \SI{1}{\atm}).
    \item Calcule el n\'umero de mol\'eculas de \ce{O3}
          por \si{\cubic\metre} usando el procedimiento del
          \textbf{Ejemplo~\ref{ej:moleculasl}} del cap\'itulo.
    \item Compare la densidad num\'erica total del aire
          a \SI{25}{\celsius} y \SI{0.77}{\atm} con la de
          \SI{25}{\celsius} y \SI{1}{\atm}.
\end{enumerate}

\textit{[Respuesta parcial:
(a) $0.120\times48/24.45 = \SI{0.235}{\milli\gram\per\cubic\metre}$;
(b) n\'umero total de mol\'eculas/m$^3$
    $= N_A P/(RT) = 6.022\times10^{23}\times0.77
    /(0.082057\times298.15\times10^{-3})
    \approx 1.90\times10^{25}$\,mol\'ec/m$^3$;
    mol\'ec de \ce{O3} = $0.120\times10^{-6}
    \times1.90\times10^{25}
    \approx 2.28\times10^{18}$\,mol\'ec/m$^3$]}

% ----- Ejercicio A8: Conversión de unidades de concentración ----------------
\item Realice las siguientes conversiones de unidades empleando
las tablas del cap\'itulo:

\begin{enumerate}[label=\alph*)]
    \item Convierta \SI{585}{\mmHg} a pascales (Pa) y a atm\'osferas.
    \item Convierta \SI{23.9}{\calory} a electronvoltios (eV).
    \item Una persona disipa energ\'ia a \SI{110}{\watt} para
          mantenerse viva. ?`Cu\'antas kilocalor\'ias por d\'ia
          necesita consumir?
    \item Convierta \SI{0.090}{ppm_v} de ozono a
          \si{\milli\gram\per\cubic\metre} a \SI{25}{\celsius}
          y \SI{1}{\atm}.
          (Peso molecular de \ce{O3} = \SI{48}{\gram\per\mole})
    \item Convierta \SI{50}{\micro\gram\per\cubic\metre} de
          \ce{SO2} a ppm$_v$ a \SI{25}{\celsius} y \SI{1}{\atm}.
          (Peso molecular de \ce{SO2} = \SI{64}{\gram\per\mole})
\end{enumerate}

\textit{[Respuesta:
(a) $585\times133.32 = \SI{77990}{\pascal}$;
    $585/760 = \SI{0.770}{\atm}$;
(b) $23.9\times4.185\times6.242\times10^{18}
\approx 6.25\times10^{19}$\,eV;
(c) $110\,\text{W}\times86400\,\text{s}/d
    \times10^{-3}\,\text{kJ/J}/4.184
    \approx \SI{2272}{\kilo\calory\per\day}$;
(d) $0.090\times48/24.45 = \SI{0.177}{\milli\gram\per\cubic\metre}$;
(e) $50\times24.45/64 = \SI{19.1}{ppb_v} = \SI{0.019}{ppm_v}$]}

% ----- Ejercicio A9: Gas de Van der Waals (avanzado) ------------------------
\item \textbf{[Avanzado]} La ecuaci\'on de van der Waals para un
gas real es:
\[
    \left(P + \frac{a}{V_m^2}\right)(V_m - b) = RT
\]
Para el \ce{CO2}: $a = \SI{3.640}{\atm\liter^2\per\mole^2}$,
$b = \SI{0.04267}{\liter\per\mole}$
(\textbf{Cuadro~\ref{tab:1}}).

\begin{enumerate}[label=\alph*)]
    \item Calcule la presi\'on de \SI{1}{\mole} de \ce{CO2}
          a \SI{300}{\kelvin} en un recipiente de \SI{1}{\liter}
          usando tanto la ley del gas ideal como la ecuaci\'on
          de van der Waals. Compare los resultados.
    \item ?`En qu\'e condiciones (alta o baja presi\'on,
          alta o baja temperatura) se espera que la diferencia
          entre ambas ecuaciones sea mayor? Justifique en
          t\'erminos de los par\'ametros $a$ y $b$.
    \item Calcule el factor de compresibilidad
          $z = PV_m/(RT)$ para el caso del inciso (a)
          seg\'un la ecuaci\'on de van der Waals.
          ?`Es el gas m\'as o menos compresible que un gas ideal?
\end{enumerate}

\textit{[Respuesta parcial:
(a) Gas ideal: $P = RT/V_m = 0.082057\times300/1
= \SI{24.62}{\atm}$;
Van der Waals: $P = RT/(V_m-b) - a/V_m^2
= 0.082057\times300/0.95733 - 3.640/1
= 25.71 - 3.64 = \SI{22.07}{\atm}$]}

\end{ejercicios}

\end{bloque_ejercicios}
% =============================================================================
% FIN DE ejercicios_cap01a.tex
% =============================================================================